%! Author = chtompki
%! Date = 1/14/26
% Example LaTeX document for GP111 - note % sign indicates a comment
\documentclass[12pt]{amsart}
% Default margins are too wide all the way around. I reset them here
\setlength{\hoffset}{-.75in}
\setlength{\textwidth}{6.5in}
\setlength{\voffset}{-.5in}
\setlength{\textheight}{9.0in}
\usepackage{amsmath}
\usepackage{amssymb}
\usepackage{amsthm}
\usepackage{pgfplots}
\usepackage{enumerate}
\usepackage{tikz}
\usetikzlibrary{arrows.meta} % Required for the Circle and Triangle arrow tips
\usepackage{setspace}

\theoremstyle{definition}
\newtheorem{definition}{Definition}

\title{Seminar Summary:\\February 3, 2026 - Allison Moore\\Cutting Knots Down to Size}
\author{Rob Tompkins}
\date{\today}


\begin{document}
    \maketitle
    \date{\today}
    \begin{onehalfspace}
        The goal behind the talk was to inspect tangle decomposition as a tool for understanding 3-manifold theory
        and $\mathbb{R}^3$ topology.
        We begin by considering a few definitions.

        \begin{definition}\;
            \begin{enumerate}[i]
                \item A \emph{knot} is an embedding of a circle into $\mathbb{R}^3$ space or an embedding into a 3-sphere.
                \item A \emph{link} $L$ is union of knots.
                \item A \emph{link diagram} is a projection of $L$ into the plane with over-crossing and under-crossings
                      explicitly labeled.
                \item \emph{Spacial graphs} or a \emph{spacial representation} $\mathcal{R}(G)$, of a graph $G$, is the
                      embedded image of $G$ in $\mathbb{R}^3$ where the vertices of $G$ are distinct points in
                      $\mathbb{R}^3$ and the edges between them are simple Jordan curves between them in such a way that
                      any two curves either are disjoint or they meet in a common end point.
                \item A \emph{theta curve} (or $\theta$-curve) is a spacial embedding of the theta-graph in
                      $\mathbb{R}^3$ (specifically in the 3-sphere $S^3$ or Euclidean space $\mathbb{R}^3$).
                \item \emph{Ambient Isotopy} defines when two knots are considered equivalent by determining if one can
                      be continuously deformed into the other within $\mathbb{R}^3$ without cutting, breaking, or
                      passing the string through itself.
            \end{enumerate}
        \end{definition}


    \end{onehalfspace}
\end{document}
